\usepackage[english]{babel}
\usepackage{color}
\usepackage{graphicx}

\usepackage{amsthm,amsmath, amsfonts}
\usepackage{amssymb,verbatim}

% algo
\usepackage{algorithm,algpseudocode}

% colors
\usepackage{color}
\usepackage[dvipsnames]{xcolor}

\definecolor{mygray}{gray}{0.1}
% Geometry of the paper
\usepackage[a4paper, margin=2cm]{geometry}

% % Appendix and bibliography
\usepackage[toc,page]{appendix}
\usepackage{authblk}

% Hyperlinks
\definecolor{DarkMidnightBlue}{rgb}{0.0, 0.2, 0.5}
\usepackage[colorlinks=true,linkcolor=Maroon,citecolor=DarkMidnightBlue,urlcolor=black]{hyperref}

% Line numbers
\usepackage{lineno}


% drawing
\usepackage{pgfplots}
\usetikzlibrary{patterns}
\usetikzlibrary{math}
\usepgfplotslibrary{fillbetween}
%<------ Or use this one

 \newcommand{\independent}{\mathop{{\rotatebox[origin=c]{90}{$\models$}}}}
\newcommand{\pder}[2][]{\frac{\partial #1}{\partial #2}}
\makeatletter
\def\lstAZ{A, B, C, D, E, F, G, H, I, J, K, L, M, N, O, P, Q, R, S, T, U, V, W, X, Y, Z}
\def\lstaz{a, b, c, d, e, f, g, h, i, j, k, l, m, n, o, p, q, r, s, t, u, v, w, x, y, z}

\def\lstAZBB{B, C, D, E, F, G, H, I, J, K, L, M, N, O, P, Q, R, T, U, V, W, X, Y, Z}

\newcommand{\MkScr}[1]{\expandafter\def\csname s#1\endcsname{\mathscr{#1}}}
\newcommand{\MkUp}[1]{\expandafter\def\csname u#1\endcsname{\mathrm{#1}}}
\newcommand{\MkFrak}[1]{\expandafter\def\csname f#1\endcsname{\mathfrak{#1}}}
\newcommand{\MkCal}[1]{\expandafter\def\csname c#1\endcsname{\mathcal{#1}}}
\newcommand{\MkBB}[1]{\expandafter\def\csname #1#1\endcsname{\mathbb{#1}}}

\@for\i:=\lstAZ\do{%
	\expandafter\MkScr \i  %  
	\expandafter\MkFrak \i  %
	\expandafter\MkUp \i %
	\expandafter\MkCal \i  %
		  }    
\@for\i:=\lstaz\do{%
	\expandafter\MkUp \i   }    
	
\@for\i:=\lstAZBB\do{%
	\expandafter\MkBB \i     }
	    
\makeatother

\newcommand{\rr}{\mathrm{r}}
\newcommand{\bb}{\mathrm{b}}
\newcommand{\zz}{\mathrm{zz}}
\newcommand{\bps}{\mathrm{bps}}
\newcommand{\Vol}{\mathrm{Vol}}

\renewcommand{\P}{\mathrm{P}}
\newcommand{\E}{\mathrm{E}}
\newcommand{\e}{\mathrm{e}}

\newcommand{\nm}[1]{\mathoperator{#1}}

% Probability
\newcommand{\PR}{\mathbb{P}}


%\newcommand{\bONE}{\mathbbm{1}}

\newcommand{\ind}{\mathbf{1}}

\newcommand{\ds}{\frac{\partial}{\partial s}}
\newcommand{\dt}{\frac{\partial}{\partial t}}
\newcommand{\dr}{\frac{\partial}{\partial r}}

\newcommand{\dd}{ \mathrm{d}}
\newcommand{\re}{ \mathrm{Re}\,}
\newcommand{\Ge}{ \mathcal{L}}
\newcommand{\sign}{\operatorname{sign}}
\DeclareMathOperator{\sgn}{sgn} 
\DeclareMathOperator*{\argmax}{argmax}
\DeclareMathOperator*{\argmin}{argmin}

\renewcommand{\epsilon}{\varepsilon}
 
\newcommand{\sfd}{{\sf d}}


\newcommand{\vn}[1]{\left| \! \left| #1\right| \! \right|} 
\newcommand{\tn}[1]{\left| \! \left| \! \left|#1\right| \! \right| \! \right|}

\newcommand{\ip}[2]{\langle #1,#2\rangle}

%moritz' definitions

\renewcommand{\tilde}{\widetilde}

\numberwithin{equation}{section}

%\theoremstyle{margin}
%\theoremseparator{.}
\newtheorem{theorem}{Theorem}[section]
\newtheorem{lemma}[theorem]{Lemma}
\newtheorem{proposition}[theorem]{Proposition}
\newtheorem{corollary}[theorem]{Corollary}

%\theoremstyle{margin}
%\theorembodyfont{\normalfont}
% \theoremstyle{definition} 
%\theoremseparator{.}
\newtheorem{definition}[theorem]{Definition}
\newtheorem{remark}[theorem]{Remark}
\newtheorem{conjecture}[theorem]{Conjecture}
\newtheorem{assumption}[theorem]{Assumption}
\newtheorem{example}[theorem]{Example}
\newtheorem{condition}[theorem]{Condition}


% short arrow
\usepackage{stmaryrd}


%---------------------------------------------------------------------------------------


\newcommand{\C}{C}

% \newcommand{\E}{\mathbb E}
\newcommand{\N}{\mathbb N}
% \renewcommand{\P}{\mathbb P}
\newcommand{\R}{\mathbb R}
\newcommand{\1}{\mathbbm 1}
\newcommand{\todo}[1]{\textcolor{blue}{\textbf{TO DO}: #1}}


% ORCID
\definecolor{lime}{HTML}{A6CE39}
\DeclareRobustCommand{\orcidicon}{%
	\begin{tikzpicture}
	\draw[lime, fill=lime] (0,0) 
	circle [radius=0.16] 
	node[white] {{\fontfamily{qag}\selectfont \tiny ID}};
	\draw[white, fill=white] (-0.0625,0.095) 
	circle [radius=0.007];
	\end{tikzpicture}
	\hspace{-2mm}
}

\foreach \x in {A, ..., Z}{%
	\expandafter\xdef\csname orcid\x\endcsname{\noexpand\href{https://orcid.org/\csname orcidauthor\x\endcsname}{\noexpand\orcidicon}}
}
% orcid
\newcommand{\orcidauthorA}{0000-0003-0185-5804}
\newcommand{\orcidauthorB}{0000-0002-2026-261X}
\newcommand{\orcidauthorC}{0000-0001-7246-8612}
\newcommand{\orcidauthorD}{0000-0003-3310-7915}
